%% ma: L11-B17-0012
%% lop: 11
%% chuong: 5
%% bai: 17
%% dang: 11.17.2
%% loai: TL
%% muc: VD
%% dap_an: "-"
%% nguon: "tai-lieu-hoc-tap-mon-toan-lop-11-theo-chuong-trinh-giao-duc-pho-thong-2018.pdf (Lớp Toán Thầy Duy, Phạm Lê Duy, THPT Chu Văn An), Chương 3 Bài 3, Dạng 2, Ví dụ 12 (tr.390-392)"
%% kien_thuc: "Lớp 11 Bài 17 (hàm số liên tục)"
%% trang_thai: cach-ly
%% ly_do: "Với đề như in thì không tồn tại a (đáp án rỗng, bất thường cho một ví dụ minh họa); nghi mẫu số của nhánh x ≠ 1 in nhầm (đúng ra là x - 1, như TLN Câu 4). Cần giáo viên xác nhận đề."
%% ghi_chu: "Đã kiểm: a ≠ -3 thì giới hạn = 0 ≠ 3 + a; a = -3 thì giới hạn = 1 ≠ 0."
\begin{ex}
Cho hàm số $f(x)=\begin{cases}\dfrac{x^3-x^2+2x-2}{3x+a}&\text{khi }x\ne1\\3x+a&\text{khi }x=1.\end{cases}$ Tìm các giá trị của tham số $a$ để $f(x)$ liên tục tại $x=1$.
\loigiai{\sloppy Với đề như in: $x^3-x^2+2x-2=(x-1)(x^2+2)$. Nếu $a\ne-3$ thì $\lim\limits_{x\to1}f(x)=0$ còn $f(1)=3+a\ne0$; nếu $a=-3$ thì $\lim\limits_{x\to1}f(x)=\dfrac{x^2+2}3=1\ne f(1)=0$. Vậy không có $a$ nào thỏa mãn.}
\end{ex}
